\chapter{Dogfooding PebbleIR}
\label{ch:pebble-dogfooding}

An IR is being \term{dogfooded} when a client uses it to solve a real implementation problem and its design changes in response to that use.
For this specimen, the client is an expression compiler that has no privileged access to optimizer internals.
Graph fixtures add the branches and loops that its source language cannot yet express.
Together they exercise the representation as a shared interface rather than as a collection of isolated structs.

\section{The client path}
\label{sec:pebble-client-path}
The source grammar accepts \code{x}, decimal literals from 0 through 255, parentheses, addition, and multiplication.
Multiplication binds more tightly; operators at the same precedence associate to the left.
The parser rejects trailing tokens, missing operands, missing closing parentheses, oversized literals, excessive nesting, and capacity exhaustion.
It builds a temporary function and publishes only a fully verified result.

The demo compiles an expression, optimizes it, lowers it, prints the optimized IR, and executes the target at input 9.
Listing~\ref{lst:pebble-client} includes the real client implementation.
The IR goes to standard output and the evaluation report to standard error, so saving the snapshot does not mix the record format with commentary.

\begin{listing}[htbp]
\inputminted{c}{example-code/pebbleir/demo.c}
\caption{A client uses the same public compile, optimize, lower, print, and execute interfaces as the tests.}
\label{lst:pebble-client}
\end{listing}

The expression \code{(x+7)*(3+5)} reports result 128 and six target dispatches.
Its unoptimized lowering has eight instructions.
The reduction is a measured static instruction-count difference checked by the test program, not a claim about hardware speedup.
The serialized value count stays seven because identity slots are retained.

\section{A testing matrix with different oracles}
\label{sec:pebble-test-matrix}
The executable suite has three graph fixtures: a diamond, a modular summation loop, and a loop-carried parallel swap.
Each runs for every input byte through the original IR, optimized IR, original VM, and optimized VM, with a direct result oracle.
It also generates 128 bounded expression trees, prints each as source, compiles it, and compares source-tree evaluation with both IR forms and optimized VM execution on every byte input.
These are 32,768 generated-program/input pairs, plus the graph cases.
They are not exhaustive over tree shapes or graphs.

Malformed-object tests target unavailable sibling definitions, incorrect phi predecessor coverage, nonboolean branch operands, out-of-range constants and targets, and invalid value-owner blocks.
Failure-contract tests exercise rejected replacement, source parsing, text parsing, value capacity, and revision exhaustion.
Round-trip tests compare printed bytes before and after parsing, including optimized tombstones.
Separate target tests check a malformed register and the five-versus-six-dispatch budget boundary.

Add a regression at the layer where the failure first becomes visible.
A parser crash needs a source reproducer; a dominance defect needs a graph reproducer; a wrong phi copy needs the edge bundle and an input that distinguishes sequential from parallel behavior.
A test that merely confirms \enquote{the program did not crash} would miss most of these failures.

\section{Reduce one failure without changing its meaning}
\label{sec:pebble-reducer}
Algorithm~\ref{alg:pebble-reduce} is a proposed reducer workflow, not an implemented feature of the library.
It keeps the same wrong-code signature, not just any failure.
Replacing a well-formed wrong-result case with malformed input that the verifier rejects does not preserve the original bug.

\begin{algorithm}[htbp]
\caption{Reducing a wrong-result reproducer while retaining its contract.}
\label{alg:pebble-reduce}
\begin{pseudocode}
\Require valid failing specimen $T$, fixed input, deterministic same-failure predicate
\Ensure failure retained and no accepted single simplifying edit remains
\Repeat
\State $changed\gets false$
\ForAll{permitted smaller candidates $T'$ in deterministic order}
\If{$T'$ verifies, all executions complete, and the same result mismatch remains}
\State $T\gets T'$; $changed\gets true$; restart candidate enumeration
\State \Break
\EndIf
\EndFor
\Until{not $changed$}
\State retain the final snapshot, input, and tool configuration
\end{pseudocode}
\end{algorithm}

For a hypothetical multiplication-lowered-as-addition defect, the expression \code{2*3} is already a useful witness: source result 6 differs from target result 5.
Replacing it by literal 6 removes the faulty multiplication and therefore removes this signature.
Keeping the operator while shrinking its operands may retain a mismatch; for example, \code{0*1} yields 0 versus 1.
This demonstrates why a reducer needs an explicit size measure and allowed edits before it claims minimality.

\section{Measure the implementation being proposed}
\label{sec:pebble-measure}
Before changing the containers, record live values, allocated slots, block count, verifier work, and target instruction count on representative clients.
Before claiming elapsed-time improvement, run repeated measurements outside sanitizer builds, record compiler and machine configuration, separate translation from execution, and verify every measured result.
The supplied suite reports correctness and the eight-to-six instruction change; it does not publish timing measurements.

Useful next clients are a statement parser that emits the existing loop graphs, a module owner that compiles several functions, and a larger graph workload that stresses repeated scans.
Each forces a specific design decision: SSA construction, cross-function identity, or indexing and allocation policy.
Adopt an external memory or container library only after specifying what new contract it must satisfy.
The point of dogfooding is to turn such decisions into observed requirements rather than preferences defended in advance.

\section{Exercises}
\label{sec:pebble-sixteen-exercises}
\begin{enumerate}
\item \label{ex:pebble-precedence} Predict the results and optimized live-value counts for \code{2+3*4} and \code{(2+3)*4}. Explain how a precedence bug could be found by the tree oracle.
\item \label{ex:pebble-reduction-signature} For the hypothetical multiply-as-add defect, evaluate \code{2*3}, \code{0*1}, and literal 6 under the correct and faulty implementations. Which candidates preserve the wrong-result signature?
\end{enumerate}

\section{Further reading}
\label{sec:pebble-sixteen-reading}
Safonov gives context for assembling multiple forms of correctness evidence around a compiler~\cite{safonov-2010-trustworthy}.
Appel remains useful as an example of learning representation requirements from the needs of a complete language implementation~\cite{appel-1998-modern}.
